% =====================================================================
%  5. Salvaguardas ambientais e sociais
% =====================================================================
\section{Salvaguardas ambientais e sociais}
\label{sec:salvaguardas}

A execução do Subcomponente 3.2.b observará o Quadro Ambiental e Social do
Banco Mundial \cite{bm-esf} por meio dos instrumentos do PSH-PR: o Marco de
Gestão Ambiental e Social (MGAS), o Plano de Engajamento das Partes
Interessadas (PEPI), os Procedimentos de Gestão de Mão de Obra (PGR), o Marco
de Política de Reassentamento (MRF) e o Mecanismo de Queixas e Reclamações
(MQR). O \autoref{qua:salvaguardas} relaciona, para cada grupo de atividades,
os riscos potenciais, as medidas de mitigação, os instrumentos aplicáveis e
o responsável pelo acompanhamento.

\begin{quadro}[htbp]
\caption{Riscos ambientais e sociais e medidas de mitigação do Subcomponente 3.2.b}
\label{qua:salvaguardas}
\footnotesize
\renewcommand{\arraystretch}{1.25}
\begin{tabularx}{\textwidth}{|P{2.4cm}|L|L|P{2.5cm}|P{1.9cm}|}
\hline
\cab{Atividade} & \cab{Riscos potenciais} & \cab{Medidas de mitigação} &
\cab{Instrumento} & \cab{Acompanha\-mento} \\ \hline
Planejamento e gestão &
Exclusão de grupos vulneráveis ou comunidades tradicionais nos diagnósticos rurais &
Consultas públicas e metodologias participativas para seleção das áreas &
PEPI & IAT \\ \hline
Sistemas de abastecimento e distribuição &
Contaminação de aquíferos; impactos temporários de obra (ruído, poeira e erosão local) &
Adoção de normas técnicas de perfuração &
MGAS & IAT \\ \hline
Operação e manutenção &
Falta de sustentabilidade técnica; conflitos locais pelo uso da água &
Treinamento das associações locais; definição de modelos de gestão e tarifa social &
MGAS e Plano de Gestão Social & IAT e prefeituras \\ \hline
Sistemas de tratamento de esgoto &
Manejo inadequado de lodo; contaminação do solo por instalação incorreta &
Capacitação das Unidades de Produção Familiar (UPF); observância da RDC 49/2013 &
MGAS e protocolos de assistência técnica & IAT / IDR-Paraná \\ \hline
Relações de trabalho nas obras &
Acidentes de trabalho ou condições inadequadas para operários e técnicos &
Protocolos de saúde e segurança do trabalho &
PGR & IAT \\ \hline
\end{tabularx}
\fonte{Adaptado de \citeonline[Quadro 26]{mop2026}.}
\end{quadro}

Aplicam-se ainda as medidas gerais do Subcomponente 3.2 \cite[Quadro 25]{mop2026}:
verificação prévia da titularidade das áreas de captação e reservação, com
anuências e termos de servidão (MRF); testes de estanqueidade após a
instalação dos sistemas de esgotamento; busca ativa de famílias em áreas
remotas e reuniões comunitárias em linguagem simples (PEPI); e implantação,
divulgação e monitoramento semestral do Mecanismo de Queixas e Reclamações
(MQR).

\pendencia{Conferir a citação da RDC 49/2013 (Anvisa). Essa resolução trata
da regularização de atividades de interesse sanitário do microempreendedor
individual e do empreendimento familiar rural, e não do manejo de lodo de
sistemas de esgotamento. Verificar qual norma se pretende citar.}
